\RequirePackage{pdfmanagement}
\SetKeys[document/metadata]{pdfversion=1.7, pdfstandard=A-2u, lang=en-GB}
\input{preamble/pdfA-compliance/glyphtounicode}
\documentclass[11pt,a4paper]{report}

\newcommand*{\PaperTitle}{Absolute concentration robustness with multistability in bimolecular networks: three species suffice, dimension two does not}
\newcommand*{\PaperTitlePlain}{Absolute concentration robustness with multistability in bimolecular networks: three species suffice, dimension two does not}
\newcommand*{\PaperShortTitle}{ACR with multistability}
\newcommand*{\PaperSubtitle}{chemistry-open-problems \textperiodcentered{} 27 October 2026 \textperiodcentered{} verified}
\newcommand*{\PaperCredits}{Solver: Claude Opus 5.5 (effort xhigh) \textperiodcentered{} Verifier: Claude Opus 5.5 (effort xhigh)}
\newcommand*{\PaperKeywords}{absolute concentration robustness, multistability, reaction networks, mass-action kinetics, bimolecular networks, deficiency}
\input{chemop-preamble}
\addbibresource{references.bib}

\providecommand{\R}{\mathbb{R}}
\newcommand{\Rpos}{\mathbb{R}_{>0}}
\newcommand{\Rnn}{\mathbb{R}_{\ge 0}}
\newcommand{\spn}{\operatorname{span}}
\newcommand{\supp}{\operatorname{supp}}
\newcommand{\one}{\mathbf{1}}
\newcommand{\Vp}{V^{+}}

\begin{document}
\makepapertitle

\begin{paperabstract}
Kaihnsa, Nguyen and Shiu proved that an at-most-bimolecular reaction network in which absolute concentration
robustness (ACR) and nondegenerate multistationarity coexist has at least 3 species, 3 reactant complexes,
5 complexes and a two-dimensional stoichiometric subspace, all four bounds being attained. They conjectured that
each of these four lower bounds becomes strictly larger when multistationarity is replaced by multistability. We
show that the conjecture holds for three of the four quantities and fails for the fourth. The main result is that a
two-dimensional bimolecular mass-action system with ACR has at most one exponentially stable positive steady state
in each stoichiometric compatibility class. The proof combines a normal form of the vector field on the ACR
hyperplane, a connectedness property of the positive steady-state set that holds precisely because products are
at most bimolecular, and the global ACR condition. Consequently ACR with multistability needs dimension at least 3,
at least 4 reactant complexes and, by a kinetic-subspace argument and the deficiency-one theorem, at least
6 complexes. On the other hand, three species suffice: the six-reaction network
$\{Z\to Y+Z,\ 2Z\to Z,\ Y\to Y+Z,\ Y+Z\to 2X,\ X+Z\to 2Z,\ 2X\to Y\}$ has ACR in $X$ for all rate constants and,
for suitable ones, two exponentially stable positive steady states, certified in exact arithmetic. It has
$3$ species, $6$ complexes and dimension $3$, so these three minima are attained. The minimal number of reactant
complexes lies between $4$ and $6$; with exactly $4$ it is impossible for $4$ species and for six of the nine
reactant structures that can occur, and with $3$ species at least $6$ reactions are needed.
\end{paperabstract}

\section*{The question}

Kaihnsa, Nguyen and Shiu~\autocite{KNS2024} study the minimal size of reaction networks in which ACR and
multistationarity coexist. Their main result (Theorem~1.2 there) states that an at-most-bimolecular network~$G$
admitting a vector of rate constants for which the mass-action system has ACR and is nondegenerately
multistationary has at least 3 species, at least 3 reactant complexes, at least 5 complexes and dimension at least
2; the network $\{X_1+X_2\to 2X_3,\ X_3\to X_1,\ 2X_3\to 2X_2\}$ attains all four bounds. Section~6 of the paper
ends with the following conjecture:

\begin{openquote}
``In the present work, our interest in multistationarity comes from the fact that it is a necessary condition for
multistability. Another interesting direction, therefore, is to investigate conditions for coexistence of ACR and
multistability, rather than multistationarity. The "minimal" networks in the current work admit only two positive
steady states and are not multistable. Hence, we conjecture that the lower bounds (on dimension and the numbers of
species, reactant complexes, and all complexes) for the coexistence of ACR and multistability are strictly larger
than the bounds proven here for ACR and multistationarity.''
\end{openquote}

\begin{theorem}[main result]\label{thm:main}
Let $G$ be an at-most-bimolecular reaction network and suppose that for some vector of positive rate constants
$\kappa$ the mass-action system $(G,\kappa)$ has ACR and is multistable. Then
\begin{enumerate}
\item the stoichiometric subspace of $G$ has dimension at least $3$;
\item $G$ has at least $4$ reactant complexes;
\item $G$ has at least $6$ complexes.
\end{enumerate}
On the other hand, there is a bimolecular network with $3$ species, $6$ complexes and dimension $3$ that has ACR for
every vector of rate constants and, for an open set of them, two exponentially stable positive steady states in one
compatibility class (\zcref{thm:three-species}). Hence the conjecture holds for dimension, reactant complexes and
complexes, and fails for the number of species. The bounds $3$ (dimension), $3$ (species) and $6$ (complexes) are
attained; the minimal number of reactant complexes lies between $4$ and $6$ (\zcref{sec:further}).
\end{theorem}

\begin{table}[ht]
\centering
\begin{tabular}{lccc}
\toprule
& ACR + nondegenerate & ACR + multistability & ACR + multistability\\
& multistationarity \autocite{KNS2024} & lower bound & best example\\
\midrule
dimension & $2$ & $3$ & $3$ ($\mathcal N_1$)\\
species & $3$ & $3$ & $3$ ($\mathcal N_1$)\\
reactant complexes & $3$ & $4$ & $6$ ($\mathcal N_1$)\\
complexes & $5$ & $6$ & $6$ ($\mathcal N_1$)\\
\bottomrule
\end{tabular}
\caption{Minimal sizes of at-most-bimolecular networks. All bounds of the first column are attained by the network
$\{X_1+X_2\to2X_3,\ X_3\to X_1,\ 2X_3\to2X_2\}$ of \autocite{KNS2024}.}
\label{tab:summary}
\end{table}

Part~1 is the substance of the paper (\zcref{thm:dim-two}); parts~2 and~3 follow from it together with results of
\autocite{KNS2024} and a short argument on the kinetic subspace (\zcref{sec:consequences}).

\section{Setting}\label{sec:setting}

A reaction network $G$ on species $X_1,\dots,X_n$ is a finite set of reactions $y\to y'$ between complexes
$y,y'\in\mathbb{Z}_{\ge0}^n$, $y\neq y'$. It is \emph{at most bimolecular} (bimolecular for short) if every
complex, reactant or product, satisfies $|y|:=y_1+\dots+y_n\le 2$. Given positive rate constants
$\kappa=(\kappa_{y\to y'})$, the mass-action system is
\begin{equation}\label{eq:ma}
\dot x=f(x):=\sum_{y\to y'}\kappa_{y\to y'}\,x^{y}\,(y'-y)=\sum_{y}x^{y}\,v_{y},\qquad
v_{y}:=\sum_{y'\colon y\to y'}\kappa_{y\to y'}(y'-y),
\end{equation}
the second sum running over reactant complexes. The stoichiometric subspace is
$S=\spn\{y'-y\}$, $s=\dim S$ is the dimension of $G$, and the stoichiometric compatibility classes are the sets
$(x^0+S)\cap\Rnn^n$. A positive steady state $x^*$ is \emph{nondegenerate} if the restriction $Df(x^*)|_S\colon S\to
S$ is invertible, and \emph{exponentially stable} if all eigenvalues of $Df(x^*)|_S$ have negative real part; the
latter implies the former.

\begin{definition}\label{def:acr-ms}
Following \autocite[Definition~2.6]{KNS2024}, $(G,\kappa)$ has \emph{ACR in $X_i$} if it has a positive steady state
and $x_i$ takes the same value $c>0$ at all positive steady states (in all compatibility classes). We call
$(G,\kappa)$ \emph{multistable} if some compatibility class contains at least two exponentially stable positive
steady states. Both properties refer to the same $\kappa$.
\end{definition}

This is the usual notion of stability in the literature on small multistable networks
\autocite[§2.2]{JTZ2025}: a steady state is stable if it is nondegenerate and the nonzero eigenvalues of the Jacobian
have negative real parts, which is equivalent to our definition. As in \autocite[Theorem~1.2]{KNS2024}, the
question is existential: a network qualifies if one vector $\kappa$ gives both properties.

Exponentially stable steady states are nondegenerate, so a multistable system is nondegenerately multistationary and
every lower bound of \autocite{KNS2024} applies to ACR with multistability.

\begin{lemma}[coefficient cones]\label{lem:cones}
Let $G$ be bimolecular and let $v_y$ be the coefficient vector of the monomial $x^y$ in~\eqref{eq:ma}. Then
$(v_y)_i\ge 0$ for every $i\notin\supp(y)$, and $\sum_i (v_y)_i\le 0$ whenever $|y|=2$.
\end{lemma}

\begin{proof}
Each $y'-y$ with $y\to y'$ has $(y'-y)_i=y'_i\ge 0$ for $i\notin\supp(y)$, and $\sum_i(y'-y)_i=|y'|-|y|\le 0$ if
$|y|=2$, because $|y'|\le 2$. Both properties pass to nonnegative combinations.
\end{proof}

The first property is the familiar characterization of mass-action vector fields~\autocite[Lemma~2.3]{KNS2024}; the
second one is where bimolecularity of the \emph{products} enters, and it is essential below
(\zcref{rem:products}).

\section{Dimension two}\label{sec:dim-two}

Throughout this section $G$ is bimolecular with $\dim S=2$, $\kappa$ is fixed, $(G,\kappa)$ has ACR in $X_1$ with
value $c>0$, $f$ is the vector field~\eqref{eq:ma} and $H:=\{x\in\R^n\colon x_1=c\}$. All positive steady states lie
in $H$. For $x\in H$ we write $x=(c,x')$ with $x'\in\R^{n-1}$.

\begin{lemma}\label{lem:not-conserved}
If $(G,\kappa)$ has a nondegenerate positive steady state, then $e_1\notin S^\perp$.
\end{lemma}

\begin{proof}
Let $x^*$ be a nondegenerate positive steady state and let $W$ be a matrix whose $n-2$ rows form a basis of
$S^\perp$. The map $F(x)=(f(x),Wx)\in S\times\R^{n-2}$ has invertible derivative at $x^*$: if $Df(x^*)u=0$ and
$Wu=0$ then $u\in S$, hence $u=0$. By the inverse function theorem, for every $T$ near $Wx^*$ there is a positive
steady state $x(T)$ near $x^*$ with $Wx(T)=T$. If $e_1\in S^\perp$ we may take $e_1^{\mathsf T}$ as a row of $W$; the
corresponding coordinate of $T$ equals $x_1(T)$, which can then be made different from $c$, contradicting ACR.
\end{proof}

Under the hypothesis of the lemma, $S\cap e_1^\perp$ is one-dimensional; we fix a vector $\ell$ spanning it. Every
compatibility class meets $H$ in a segment parallel to~$\ell$.

\begin{lemma}[the field on the ACR hyperplane]\label{lem:field-on-H}
Suppose some compatibility class contains two distinct nondegenerate positive steady states $p$ and $q$. Then there
are $\alpha\in S\setminus\{0\}$ and a polynomial $h$ of degree at most $2$ such that $f(x)=h(x)\,\alpha$ for all
$x\in H$. Moreover $q=p+\tau\ell$ for some $\tau\neq0$, and $h_L(t):=h(p+t\ell)$ has degree exactly $2$ with simple
zeros $t=0$ and $t=\tau$.
\end{lemma}

\begin{proof}
Choose a basis $\sigma_1,\sigma_2$ of $S$ and write $f=g_1\sigma_1+g_2\sigma_2$ with polynomials $g_1,g_2$ of degree
at most $2$. Since $Df(p)$ maps onto $S$, the set $f^{-1}(0)$ is, near $p$, a smooth submanifold of dimension $n-2$,
and by ACR its positive points lie in~$H$. Thus the restrictions $g_1|_H$ and $g_2|_H$, viewed as polynomials on
$H\cong\R^{n-1}$, vanish on a smooth hypersurface $M$ of $H$.

We claim that they have a nonconstant common factor (or one of them vanishes identically, in which case the lemma
holds with $\alpha$ the other basis vector). By Baire's theorem, an irreducible factor $F_1$ of $g_1|_H$ vanishes on a
nonempty open subset $M_1$ of $M$, and an irreducible factor $F_2$ of $g_2|_H$ vanishes on a nonempty open subset
$M_2\subseteq M_1$. If $F_1$ and $F_2$ were not proportional, they would be coprime, their common complex zero set
would have dimension at most $n-3$, and so would its real points; but it contains the $(n-2)$-dimensional manifold
$M_2$.

If the common factor $F$ has degree $2$, then $g_i|_H=c_iF$ with constants $c_i$, and $f|_H=F\,(c_1\sigma_1+c_2\sigma_2)$.
If $F$ has degree $1$, then $g_i|_H=F\,m_i$ with $\deg m_i\le 1$. If $m_1,m_2$ are proportional we conclude as before.
Otherwise $\{x\in H\colon m_1=m_2=0\}$ is empty or an affine subspace of dimension at most $n-3$. A nondegenerate
positive steady state $y$ with $F(y)\neq 0$ would have a neighbourhood in $f^{-1}(0)$, of dimension $n-2$, contained
in that set; so every nondegenerate positive steady state satisfies $F=0$. The steady states of the class of $p$ lie
on the segment $(p+\R\ell)\cap\Rpos^n$, on which $F$ is affine; vanishing at $p\neq q$, it vanishes on the whole
segment, so $p$ is not isolated in its class, which contradicts nondegeneracy (nondegenerate steady states are
isolated in their class by the inverse function theorem).

Finally, $p,q\in H$ lie in the same class, so $q-p\in S\cap e_1^\perp=\R\ell$. The polynomial $h_L$ vanishes at $0$
and $\tau$; if its degree were at most one it would vanish identically and $p$ would again not be isolated. Hence
$\deg h_L=2$ and both zeros are simple. Likewise $f|_H\not\equiv0$ (otherwise $p$ would not be isolated), so the
vector $\alpha$ obtained above is nonzero.
\end{proof}

From now on we assume that a class contains two distinct nondegenerate positive steady states $p,q$, and we use
$\alpha$, $h$, $\ell$, $\tau$ and $h_L$ from \zcref{lem:field-on-H}.

\begin{lemma}\label{lem:alpha1-zero}
If $\alpha_1=0$, then $Df|_S$ has a positive real eigenvalue at $p$ or at $q$.
\end{lemma}

\begin{proof}
Then $\alpha\in S\cap e_1^\perp$, so $\alpha=\gamma\ell$ with $\gamma\neq0$. For $x_t:=p+t\ell\in H$ we have
$f(x_t)=\gamma h_L(t)\ell$, hence $Df(x_t)\ell=\gamma h_L'(t)\ell$: $\ell$ is an eigenvector with eigenvalue
$\gamma h_L'(t)$. Since $h_L'(0)$ and $h_L'(\tau)$ are nonzero with opposite signs, this eigenvalue is positive at $p$
or at $q$.
\end{proof}

\paragraph{Normal form.}
If $\alpha_1\neq0$ we normalize $\alpha_1=1$. Since $\{\alpha,\ell\}$ is a basis of $S$ and $\ell_1=0$, we can write
uniquely
\begin{equation}\label{eq:normal-form}
f=\psi\,\alpha+\varphi\,\ell,\qquad \psi=f_1 .
\end{equation}
On $H$ we have $f=h\alpha$, so $\varphi$ vanishes on $H$ and $\psi|_H=h$. As $\deg\varphi\le2$, there is an affine
function $P$ with $\varphi=(x_1-c)\,P$. The positive steady states are the positive solutions of $\psi=0$ and
$(x_1-c)P=0$.

\begin{lemma}[determinant on the ACR hyperplane]\label{lem:det}
If $\alpha_1=1$, then $\det\bigl(Df(x)|_S\bigr)=-P(x)\,\bigl(\nabla\psi(x)\cdot\ell\bigr)$ for every $x\in H$.
\end{lemma}

\begin{proof}
On $H$, $Df=\alpha\,\nabla\psi^{\mathsf T}+P\,\ell\,e_1^{\mathsf T}$. Hence $Df\,\alpha=(\nabla\psi\cdot\alpha)\alpha+P\ell$
and $Df\,\ell=(\nabla\psi\cdot\ell)\alpha$. In the basis $(\alpha,\ell)$ the matrix of $Df|_S$ is
$\begin{psmallmatrix}\nabla\psi\cdot\alpha&\nabla\psi\cdot\ell\\ P&0\end{psmallmatrix}$, with determinant
$-P\,(\nabla\psi\cdot\ell)$.
\end{proof}

\begin{corollary}\label{cor:P-changes-sign}
If $\alpha_1=1$ and $p,q$ are both exponentially stable, then $P(p)\,P(q)<0$.
\end{corollary}

\begin{proof}
A $2\times2$ matrix whose eigenvalues have negative real parts has positive determinant. Here
$\nabla\psi(p)\cdot\ell=h_L'(0)$ and $\nabla\psi(q)\cdot\ell=h_L'(\tau)$ are nonzero with opposite signs, so
\zcref{lem:det} forces $P(p)$ and $P(q)$ to have opposite signs.
\end{proof}

The remaining steps show that $P$ cannot change sign along the positive steady states. Write
$\psi_H(x'):=\psi(c,x')$ and
\[
\Vp:=\{x'\in\Rpos^{n-1}\colon \psi_H(x')=0\},
\]
so that the positive steady states are exactly the points $(c,x')$ with $x'\in\Vp$.

\begin{lemma}[bimolecularity makes $\Vp$ connected]\label{lem:connected}
Assume $\alpha_1=1$. Then $\Vp$ is connected, and $\nabla\psi_H\neq0$ at every point of $\Vp$.
\end{lemma}

\begin{proof}
For $2\le j\le k\le n$ let $a_{jk}$ be the coefficient of $x_jx_k$ in $\psi=f_1$. Monomials containing $x_1$ become
of degree at most one in $x'$ on $H$, and $(x_1-c)P$ contains no monomial $x_jx_k$ with $j,k\ge2$; hence the
quadratic part of $\psi_H$ is $Q(x')=\sum_{j\le k}a_{jk}x_jx_k$ and the coefficient vector of $x_jx_k$ in $f$ is
$a_{jk}\alpha$. By \zcref{lem:cones} applied to $y=X_j+X_k$: $a_{jk}\ge0$ (coordinate $1$), and if $a_{jk}>0$ then
$\alpha_i\ge0$ for all $i\notin\{1,j,k\}$ and $1+\sum_{i\ge2}\alpha_i\le 0$, so that
$\sum_{i\in\{j,k\}}\alpha_i\le -1-\sum_{i\notin\{1,j,k\}}\alpha_i\le -1$ (a single term when $j=k$). Let
$N:=\{i\ge2\colon\alpha_i<0\}$. Every monomial
$x_jx_k$ with $a_{jk}>0$ therefore satisfies $N\subseteq\{j,k\}$ and $N\cap\{j,k\}\neq\emptyset$. Since
$h_L(t)=\psi_H(p'+t\ell')$ has degree $2$, $Q(\ell')\ne0$, so $Q\neq0$ and $1\le|N|\le2$.

Second, for $i\in N$ the component $f_i$ restricted to $H$ equals $\alpha_i\psi_H$; its monomials not containing
$x_i$ have nonnegative coefficients (setting $x_1=c>0$ preserves this property of $f_i$), so the monomials of
$\psi_H$ not containing $x_i$ have nonpositive coefficients. Two cases remain.

\emph{Case $N=\{a,b\}$.} Then $Q=\beta x_ax_b$ with $\beta>0$, every other monomial of $\psi_H$ lacks $x_a$ or $x_b$,
and $\psi_H=\beta x_ax_b-G_0-G_ax_a-G_bx_b-\sum_{k\notin\{a,b\}}G_kx_k$ with all $G\ge0$. Writing
$F:=G_0+G_ax_a+\sum_{k\notin\{a,b\}}G_kx_k$, the equation $\psi_H=0$ reads $x_b(\beta x_a-G_b)=F$. On the open
orthant either $F>0$, and then $\Vp$ is the graph $x_b=F/(\beta x_a-G_b)$ over the convex set
$\{\beta x_a>G_b\}\cap\Rpos^{n-2}$, with $\partial\psi_H/\partial x_b=\beta x_a-G_b>0$; or $F\equiv0$, and then
$\Vp=\{x_a=G_b/\beta\}\cap\Rpos^{n-1}$ is convex, with $\partial\psi_H/\partial x_a=\beta x_b>0$.

\emph{Case $N=\{a\}$.} Then $Q=x_a\sum_{k\ge2}\beta_kx_k$ with all $\beta_k\ge0$, not all zero, and
$\psi_H=\beta_ax_a^2+x_a\bigl(K+R(x'')\bigr)-G(x'')$, where $x''$ collects the coordinates other than $x_1,x_a$,
$R(x'')=\sum_{k\neq a}\beta_kx_k$, $K\in\R$, and $G$ is affine with nonnegative coefficients. On the open orthant
either $G>0$ or $G\equiv0$. If $G>0$ and $\beta_a>0$, the quadratic in $x_a$ has exactly one positive root, which is
simple; if $G>0$ and $\beta_a=0$, the positive root $x_a=G/(K+R)$ exists exactly on the convex set $\{K+R>0\}$, where
$\partial\psi_H/\partial x_a=K+R>0$. In both cases $\Vp$ is the graph of a smooth function of $x''$ over a convex set.
If $G\equiv0$, then $\Vp=\{\beta_ax_a+K+R(x'')=0\}\cap\Rpos^{n-1}$ is convex, and on it
$\nabla\psi_H=x_a\nabla(\beta_ax_a+K+R)\neq0$, because $\beta_a$ or some $\beta_k$ with $k\ne a$ is positive.
\end{proof}

\begin{lemma}[global ACR]\label{lem:global}
Assume $\alpha_1=1$ and $P(p)P(q)<0$, and write $P_H(x'):=P(c,x')$. Then
$Z:=\{x'\in\Rpos^{n-1}\colon P_H(x')=0\}$ is nonempty, and there is $\varepsilon\in\{1,-1\}$ with
$\varepsilon\,\psi_H\ge0$ on $Z$.
\end{lemma}

\begin{proof}
$P_H$ takes both signs on the convex set $\Rpos^{n-1}$, so $Z\neq\emptyset$; in particular $P$ depends on some
$x_k$ with $k\ge2$, and the hyperplane $\Pi:=\{P=0\}$ of $\R^n$ is transverse to $H$. At a positive point of $\Pi$
we have $f=\psi\alpha$, so by ACR $\psi$ has no zero on the convex sets
$\Pi^{\pm}:=\Pi\cap\Rpos^n\cap\{\pm(x_1-c)>0\}$, and hence has a constant sign $\varepsilon_\pm$ on each. Both sets
accumulate at every point of $\{c\}\times Z$ (move inside $\Pi$ along $e_1-(\pi_1/\pi_k)e_k$, where $\pi_j$ are the
coefficients of $P$ and $\pi_k\neq0$), so $\varepsilon_+\psi_H\ge0$ and $\varepsilon_-\psi_H\ge0$ on $Z$. If
$\varepsilon_+\neq\varepsilon_-$, then $\psi_H$ vanishes on $Z$, a nonempty open subset of the hyperplane
$\{P_H=0\}$, so $P_H$ divides $\psi_H$: $\psi_H=P_H\,A$ with $A$ affine. Since $P_H(p'),P_H(q')\neq0$, $A$ vanishes at
$p'$ and $q'$, hence on the line through them, and $h_L\equiv0$, contradicting \zcref{lem:field-on-H}. So
$\varepsilon:=\varepsilon_+=\varepsilon_-$ works.
\end{proof}

\begin{lemma}[no transversal sign change]\label{lem:tangency}
In the situation of \zcref{lem:global}, let $z\in\Vp\cap Z$. Then $\nabla\psi_H(z)=\mu\,\nabla P_H$ for some
$\mu\neq0$, and $z$ has a neighbourhood $U$ in $\Vp$ on which $-\operatorname{sign}(\mu)\,\varepsilon\,P_H\ge0$.
\end{lemma}

\begin{proof}
Let $E:=\{P_H=0\}$, with direction space $T$. The function $\varepsilon\psi_H$ is nonnegative on the relatively open
set $Z\subseteq E$ and vanishes at $z$, so $\nabla\psi_H(z)\perp T$, that is $\nabla\psi_H(z)=\mu\nabla P_H$; and
$\mu\neq0$ by \zcref{lem:connected}. Let $Q$ be the quadratic part of $\psi_H$ and $B$ its polar form. For $v\in T$,
$\psi_H(z+v)=Q(v)$, so $\varepsilon Q\ge0$ on $T$. Fix $\nu$ with $\nabla P_H\cdot\nu=1$ and write $w=v+s\nu$ with
$v\in T$, $s=P_H(z+w)$. If $z+w\in\Vp$,
\[
0=\psi_H(z+w)=\mu s+Q(v)+2sB(v,\nu)+s^2Q(\nu),\quad\text{i.e.}\quad s\,\bigl(\mu+2B(v,\nu)+sQ(\nu)\bigr)=-Q(v).
\]
For $w$ small the bracket has the sign of $\mu$, so $s$ is zero or has the sign of $-\mu\varepsilon$.
\end{proof}

\begin{theorem}[dimension two]\label{thm:dim-two}
Let $G$ be an at-most-bimolecular network of dimension $2$ and let $\kappa$ be such that $(G,\kappa)$ has ACR. Then
every stoichiometric compatibility class contains at most one exponentially stable positive steady state. In
particular $(G,\kappa)$ is not multistable.
\end{theorem}

\begin{proof}
Suppose $p\neq q$ are exponentially stable positive steady states in one class; they are nondegenerate. By
\zcref{lem:not-conserved,lem:field-on-H} we have the objects $\alpha,\ell,h$; by \zcref{lem:alpha1-zero},
$\alpha_1\neq0$, and we normalize $\alpha_1=1$. By \zcref{cor:P-changes-sign}, $P(p)P(q)<0$. Define
$\theta\colon\Vp\to\{1,-1\}$ by $\theta(x')=\operatorname{sign}P_H(x')$ if $P_H(x')\neq0$ and
$\theta(x')=-\operatorname{sign}(\mu(x'))\,\varepsilon$ if $x'\in Z$, where $\mu(x')=\nabla\psi_H(x')\cdot\nabla
P_H/|\nabla P_H|^2$ is continuous and nonzero on $\Vp\cap Z$ (\zcref{lem:tangency}). Then $\theta$ is locally
constant: near a point with $P_H\neq0$ by continuity, and near $z\in\Vp\cap Z$ by \zcref{lem:tangency} applied on a
neighbourhood where $\mu$ keeps its sign. As $\Vp$ is connected (\zcref{lem:connected}), $\theta$ is constant, so
$P(p)$ and $P(q)$ have the same sign, a contradiction.
\end{proof}

\begin{remark}[where the bimolecular products matter]\label{rem:products}
The bound $|y'|\le2$ on product complexes is used only through the second property in \zcref{lem:cones}, inside
\zcref{lem:connected}. Without it, $\psi_H$ may contain the terms $x_2^2-3x_2+2$, and $\Vp$ may consist of two
parallel sheets on which $P$ has opposite signs; the argument then breaks down. In the planar model
$\dot u=t^2-1-u/c$, $\dot t=-ut$ (with $u=x_1-c$), the two stable nodes $t=\pm1$ of the ACR line $u=0$ are separated
by the invariant line $t=0$, which ends at a saddle on the face $x_1=0$. We have not found a mass-action realization
of this mechanism with bimolecular reactant complexes (linear programs with the constraints of \zcref{lem:cones}
minus the product bound are feasible locally, but the global ACR condition failed on every sample), so we do not
know whether \zcref{thm:dim-two} extends to networks with larger product complexes.
\end{remark}

\section{Reactant complexes and complexes}\label{sec:consequences}

\begin{corollary}[dimension and reactant complexes]\label{cor:reactants}
If a bimolecular network $G$ admits $\kappa$ with ACR and multistability, then $\dim S\ge3$ and $G$ has at least
$4$ reactant complexes.
\end{corollary}

\begin{proof}
One-dimensional networks never have ACR with nondegenerate multistationarity \autocite[Proposition~3.10]{KNS2024},
and dimension~$2$ is excluded by \zcref{thm:dim-two}. A bimolecular network with ACR and nondegenerate
multistationarity has at least $3$ species \autocite[Theorem~4.1]{KNS2024}. If $G$ has $k\ge1$ independent
conservation laws, it has at least $n-k+1=\dim S+1\ge4$ reactant complexes \autocite[Theorem~5.5]{KNS2024}; if it
is full-dimensional, at least $n+2\ge5$ \autocite[Theorem~5.1]{KNS2024}.
\end{proof}

For the number of complexes we use the classical objects of chemical reaction network theory: $m$ complexes, $\ell$
linkage classes, $t$ terminal strong linkage classes, deficiency $\delta=m-\ell-s$, and the factorization
$f(x)=Y A_\kappa \Psi(x)$, where $Y$ is the $n\times m$ matrix of complexes, $\Psi(x)=(x^{y})_y$ and $A_\kappa$ is
the Laplacian matrix of the weighted reaction graph \autocite{Feinberg2019}. The kinetic subspace
$K_\kappa:=Y(\operatorname{Im}A_\kappa)\subseteq S$ contains every value of $f$ and every column of $Df$.

\begin{lemma}[kinetic subspace]\label{lem:kinetic}
If $(G,\kappa)$ has a nondegenerate positive steady state and $G$ has a nonterminal complex, then $\delta\ge
t-\ell+1$.
\end{lemma}

\begin{proof}
At a nondegenerate positive steady state $x^*$ the image of $Df(x^*)|_S$ is $S$, so $K_\kappa=S$. The kernel of
$A_\kappa$ has dimension $t$ and consists of vectors supported on terminal complexes, so
$\dim\operatorname{Im}A_\kappa=m-t$, and $A_\kappa\Psi(x^*)\neq0$ because $\Psi(x^*)$ is positive on a nonterminal
complex. As $YA_\kappa\Psi(x^*)=f(x^*)=0$, the map $Y$ has a nonzero kernel on $\operatorname{Im}A_\kappa$. Hence
$s=\dim K_\kappa\le m-t-1$, that is $\delta=m-\ell-s\ge t-\ell+1$.
\end{proof}

\begin{corollary}[complexes]\label{cor:complexes}
If a bimolecular network $G$ admits $\kappa$ with ACR and multistability, then $G$ has at least $6$ complexes.
\end{corollary}

\begin{proof}
By \zcref{cor:reactants}, $s\ge3$. A network of deficiency zero is not multistationary \autocite[Lemma~2.4]{KNS2024},
so $\delta\ge1$ and $m=\ell+s+\delta\ge5$, with equality only if $\ell=1$, $s=3$, $\delta=1$. Suppose $m=5$. If $G$
has a nonterminal complex, \zcref{lem:kinetic} gives $t\le\ell+\delta-1=1$; otherwise all complexes are terminal and
$t=\ell=1$. So $G$ has one linkage class, of deficiency one, with a single terminal strong linkage class, and the
deficiency-one theorem \autocite{Feinberg1995} gives at most one positive steady state in each compatibility class.
This contradicts multistability.
\end{proof}

\begin{remark}[on the deficiency-one theorem]\label{rem:def-one}
The uniqueness part of the deficiency-one theorem holds without weak reversibility, provided each linkage class
contains exactly one terminal strong linkage class, $\delta_\theta\le1$ for each linkage class and
$\sum_\theta\delta_\theta=\delta$ \autocite{Feinberg1995}. \textcite[Lemma~2.5]{KNS2024} state it for weakly reversible
networks only, but their proof of the bound of $5$ complexes (Theorem~4.9(2) there) applies it to networks that need
not be weakly reversible; \zcref{lem:kinetic} supplies the missing hypothesis $t=1$ there, so their bound stands. The
same argument shows that two nondegenerate positive steady states in one class force $m\ge s+3$ for networks of any
molecularity: $\delta\ge1$ always, and $\delta\ge2$ when $\ell=1$.
\end{remark}

\section{Three species and six complexes suffice}\label{sec:three}

\begin{theorem}[three species, six complexes]\label{thm:three-species}
The bimolecular network on species $X,Y,Z$
\[
\mathcal N_1:\quad Z\xrightarrow{k_1}Y+Z,\quad 2Z\xrightarrow{k_2}Z,\quad Y\xrightarrow{k_3}Y+Z,\quad
Y+Z\xrightarrow{k_4}2X,\quad X+Z\xrightarrow{k_5}2Z,\quad 2X\xrightarrow{k_6}Y
\]
has $3$ species, $6$ reactions, $6$ reactant complexes, $6$ complexes and dimension $3$. For every vector of rate
constants it has ACR in $X$, with value $2k_1/k_5$. For $(k_1,\dots,k_6)=(5,1,1,4,10,7)$ its positive steady states
are exactly
\[
P_1=\bigl(1,\tfrac{19}{4},\tfrac12\bigr),\qquad P_2=(1,3,1),\qquad P_3=\bigl(1,\tfrac74,\tfrac72\bigr),
\]
and $P_1,P_3$ are exponentially stable while $P_2$ is a saddle. The network being full-dimensional, all three lie in
the unique compatibility class $\Rnn^3$. Multistability and ACR persist for all rate constants in a neighbourhood of
this vector.
\end{theorem}

\begin{proof}
The mass-action system is
\begin{align*}
\dot x&=2k_4yz-k_5xz-2k_6x^2,\\
\dot y&=k_1z-k_4yz+k_6x^2,\\
\dot z&=-k_2z^2+k_3y-k_4yz+k_5xz .
\end{align*}
The reaction vectors $(0,1,0)$, $(0,0,-1)$, $(2,-1,-1)$ and $(-1,0,1)$ span $\R^3$. The weight $w=(1,2,0)$ gives
$\dot x+2\dot y=z\,(2k_1-k_5x)$, so every positive steady state has $x=c:=2k_1/k_5$. On this plane $\dot y=0$ gives
$y=(k_1z+k_6c^2)/(k_4z)$, and then
\[
k_4z\,\dot z=-k_2k_4z^3+k_1k_4z^2+(k_1k_3-k_4k_6c^2)\,z+k_3k_6c^2 ,
\]
which is positive at $z=0$ and negative for large $z$. So a positive steady state exists for every $\kappa$, and ACR
holds for every $\kappa$. For the given rate constants $c=1$, $y=(5z+7)/(4z)$ and
$4z\,\dot z=-(2z-1)(z-1)(2z-7)$, which gives exactly $P_1,P_2,P_3$. The Jacobian matrices and their characteristic
polynomials are
\begin{align*}
J(P_1)&=\begin{psmallmatrix}-33&4&28\\14&-2&-14\\5&-1&-10\end{psmallmatrix}, &
\chi_1&=\lambda^3+45\lambda^2+206\lambda+30=(\lambda+5)(\lambda^2+40\lambda+6),\\
J(P_2)&=\begin{psmallmatrix}-38&8&14\\14&-4&-7\\10&-3&-4\end{psmallmatrix}, &
\chi_2&=\lambda^3+46\lambda^2+47\lambda-50,\\
J(P_3)&=\begin{psmallmatrix}-63&28&4\\14&-14&-2\\35&-13&-4\end{psmallmatrix}, &
\chi_3&=\lambda^3+81\lambda^2+632\lambda+1050 .
\end{align*}
By the Routh--Hurwitz criterion ($45\cdot206>30$ and $81\cdot632>1050$, all coefficients positive) $P_1$ and $P_3$
are exponentially stable; $\chi_2(0)=-50<0$ gives $P_2$ a positive real eigenvalue. All three steady states are
nondegenerate, so they persist, with their stability types, for nearby rate constants, and ACR holds for all rate
constants.
\end{proof}

The mechanism is a weighted integral feedback: the pair of complexes $\{Z,\,X+Z\}$ is the only one whose reactions
change $w\cdot x=x+2y$, so that $w\cdot f=z\,(2k_1-k_5x)$. Two further certified networks with six complexes, one of
them weakly reversible with two linkage classes of deficiency zero and total deficiency one, are given in the
sub-study \texttt{complexes/report.md}.

\section{Further results and open questions}\label{sec:further}

The following results complete \zcref{tab:summary}. Their proofs, partly computer-assisted with exact arithmetic, are
in the two sub-study reports \texttt{examples/report.md} and \texttt{complexes/report.md}, with the scripts that
check them.

\begin{enumerate}
\item \emph{A second three-species example with rational steady states.} The network
$\mathcal G_7=\{X+Y\xrightarrow{1}2Y,\ Y\xrightarrow{1}2Z,\ 0\xrightarrow{5}X+Z,\ X+Z\xrightarrow{4}2Z,\
2Z\xrightarrow{1}X+Z,\ Y+Z\xrightarrow{6}Y\}$ has $\dot y=y(k_1x-k_2)$, hence ACR in $X$ for all rate constants,
and for the indicated ones exactly three positive steady states $(1,2,1)$, $(1,\frac54,\frac32)$,
$(1,\frac{10}9,\frac53)$, the first and the last exponentially stable (exact Gröbner basis and Routh--Hurwitz).
It has $7$ complexes.
\item \emph{Six complexes in two linkage classes.} The weakly reversible network
$\{2Z\rightleftarrows X+Z,\ Y+Z\to X,\ Y+Z\to X+Y,\ X\to2X,\ X+Y\to X,\ 2X\to Y+Z\}$ has $6$ complexes, two linkage
classes of deficiency~$0$ and total deficiency~$1$, ACR in $X$ and, for suitable rate constants, two exponentially
stable positive steady states. So the deficiency-one theorem cannot be applied to two linkage classes when
$\sum_\theta\delta_\theta<\delta$, and six-complex examples need not have a single linkage class.
\item \emph{Reactions with three species.} No full-dimensional bimolecular network with $3$ species and at most
$5$ reactions has three nondegenerate positive steady states, or two exponentially stable ones, for any rate
constants. Networks with at most $4$ distinct reactant complexes are covered by
\autocite[Propositions~3.1(3) and~3.5]{KNS2024}; the $3\,188\,646$ networks with $5$ distinct reactant complexes
were enumerated exactly, through a reduction to the sign changes of one univariate function whose derivative has,
up to a fixed sign, the sign of the Jacobian determinant (two exponentially stable states have determinants of the same sign, which
forces a third steady state between them). By \zcref{cor:reactants}, a three-species network with ACR and
multistability is full-dimensional, so it needs at least $6$ reactions, which $\mathcal N_1$ and $\mathcal G_7$
attain. The restriction to full-dimensional networks is necessary: the two-dimensional network
$\{0\xrightarrow{10}Z,\ Y+Z\xrightarrow{10}X,\ X+Z\xrightarrow{1}Y+Z,\ X+Y\xrightarrow{1}2Y\}$ has the three
nondegenerate positive steady states $(5,1,1)$, $(4,2,\frac12)$, $(3,3,\frac13)$ in the class $x+y=6$, the outer two
exponentially stable (it has no ACR).
\item \emph{Four reactant complexes.} With exactly $4$ reactant complexes and ACR, the positive steady states form a
coset of a toric variety, and the number of them in a class is bounded by a rule of signs for exponential sums.
For $4$ species no class contains two exponentially stable positive steady states (two such states would force a
third one, and the rule of signs allows at most two). For $5$ or more species exactly nine reactant
structures are compatible with ACR (none with $7$ or $8$ species); six of them are excluded by proof, and the
remaining three were excluded only numerically, by sampling all realizable stoichiometric subspaces.
\end{enumerate}

\paragraph{Open questions.} Is the minimal number of reactant complexes $4$, $5$ or $6$? For full-dimensional
networks on three species, is the bound $n+2=5$ of \autocite[Theorem~5.1]{KNS2024} attained with multistability? Does \zcref{thm:dim-two} hold for networks with
bimolecular reactant complexes but larger products (\zcref{rem:products})? Finally, \autocite{JNJ2026} found ACR with
bistability in a dual-site futile cycle with bifunctional enzymes (six species, dimension four); our examples show
that much smaller mechanisms, based on a weighted integral feedback, suffice.

\section{Computer verification}\label{sec:computer}

All scripts print \texttt{PASS}/\texttt{FAIL} for each claim and exit with a nonzero code on failure. They were run
under WSL (Fedora~44) with Python~3.12, sympy~1.14 and python-flint~0.9.
\begin{itemize}
\item \texttt{scripts/check\_dimension\_two\_example.py} (8~s): on $20$ random rational instances of the KNS network
of Example~4.10, checks the determinant formula of \zcref{lem:det} and that the two positive steady states of a
class have determinants of opposite signs. This is a sanity check; \zcref{thm:dim-two} is proved by hand.
\item \texttt{complexes/scripts/certify\_examples.py} (exact rational arithmetic, standard library only, under
$0.1$~s): \zcref{thm:three-species} and the two further six-complex examples (ACR identity, steady states,
characteristic polynomials, Routh--Hurwitz conditions, counts of complexes and deficiency).
\texttt{complexes/scripts/eigen\_enclosures.py}: eigenvalue enclosures in ball arithmetic.
\item \texttt{examples/scripts/certify\_example.py} (under $1$~s): the network $\mathcal G_7$ (exact Gröbner basis,
root isolation, ball-arithmetic eigenvalues). \texttt{examples/scripts/enum5\_reactions.py} (about two minutes on
$12$ cores): the enumeration of item~3. \texttt{examples/scripts/e2\_reactant\_sets.py} and
\texttt{check\_gale\_formula.py}: the reactant structures of item~4 and the determinant formula used there.
\end{itemize}
Searches and simulations used to find the examples are in \texttt{explore} sub-folders; they are floating-point and
not part of any proof.

\paperreferences
\end{document}
